\chapter{Vestibular Neuritis}

\section*{Definition}
Vestibular neuritis is an acute, unilateral, selective inflammation of the vestibular portion of cranial nerve VIII (vestibular nerve) causing rapid-onset vertigo, nausea, and imbalance without hearing loss or tinnitus.

\section*{Pathophysiology}
Reactivation of latent herpes simplex virus type 1 within the vestibular ganglion (Scarpa ganglion) causes inflammation and edema of the superior vestibular nerve. This produces a sudden asymmetry in tonic vestibular signals between the two sides, generating a false sense of rotation.

\section*{Etiology}
\begin{itemize}
  \item \textbf{Viral infection:} HSV-1 reactivation (primary suspect), with antecedent or concurrent URI in many cases
  \item \textbf{Other viruses:} Varicella-zoster, EBV, CMV, influenza, and SARS-CoV-2 reported associations
  \item \textbf{Immunemediated:} Post-viral inflammatory response directed at vestibular nerve antigens
  \item \textbf{Risk factors:} Stress, immunosuppression, and recent viral illness
\end{itemize}

\section*{Indications for Evaluation}
Seek care if:
\begin{itemize}
  \item Acute, severe, continuous vertigo lasting > 12 hours without hearing loss
  \item Inability to walk or stand without support due to imbalance
  \item Nystagmus that persists when visual fixation is prevented (spontaneous nystagmus in the dark)
  \item Head motion intolerance (all head movements provoke severe symptoms)
  \item Distinguish from central causes: normal neurologic exam (no dysarthria, dysmetria, or weakness)
\end{itemize}

\section*{Related Symptoms}
\begin{itemize}
  \item Severe rotational vertigo (acute phase lasts 12--72 hours)
  \item Nausea and vomiting
  \item Nystagmus (horizontal, unidirectional, fast phase toward the unaffected ear)
  \item Gait instability and postural imbalance
  \item Oscillopsia (sensation that the visual world is bouncing with head movement)
\end{itemize}
\textbf{Note:} The absence of hearing loss and tinnitus distinguishes vestibular neuritis from labyrinthitis; the absence of neurologic signs distinguishes it from central vertigo (brainstem or cerebellar stroke).

\section*{Self-Care / Home Management}
\begin{itemize}
  \item Bed rest in a quiet, dark room during the acute vertiginous phase
  \item Vestibular suppressants (meclizine, dimenhydrinate) limited to the first 1--3 days to avoid central compensation delay
  \item Avoid driving, operating machinery, and climbing
  \item Fall prevention: use a cane or walker, remove trip hazards at home
  \item Oral ginger for adjunctive antiemetic effect
\end{itemize}

\section*{Diagnostic Approach}
\begin{itemize}
  \item HINTS examination (Head Impulse, Nystagmus, Test of Skew) to differentiate peripheral from central vertigo
  \item Head impulse test: Corrective saccade toward the affected side (positive peripheral sign)
  \item Spontaneous nystagmus: Horizontal, unidirectional, suppressed by visual fixation
  \item Romberg test: Falls toward the affected side
  \item Audiometry (normal in vestibular neuritis; SNHL in labyrinthitis)
  \item MRI brain and internal auditory canals if central signs, negative HINTS, or atypical presentation
  \item Videonystagmography or caloric testing to confirm reduced vestibular function on the affected side
\end{itemize}

\section*{Treatment / Management Overview}
\begin{itemize}
  \item \textbf{Acute (first 1--3 days):} Vestibular suppressants (diazepam 2--5 mg, meclizine 25 mg, or dimenhydrinate 50 mg PO q8h) and antiemetics (ondansetron)
  \item \textbf{Short-course corticosteroids:} Methylprednisolone taper (starting 100 mg daily with rapid taper) may improve recovery if initiated within 72 hours
  \item \textbf{Vestibular rehabilitation therapy (VRT):} Habituation, gaze stabilization, and balance exercises started after the acute phase resolves
  \item \textbf{Antiviral therapy:} Acyclovir or valacyclovir of uncertain benefit; not routinely recommended
  \item \textbf{Subacute recovery:} Cessation of vestibular suppressants after day 3 to allow central compensation; gradual return to normal activity
  \item \textbf{Prognosis:} Most patients recover within weeks; persistent imbalance and motion sensitivity may require ongoing VRT
\end{itemize}

\clearpage
